\newcommand{\LeakText}{A row-wise split lowers the error by a factor of 1.2--1.8, and copying the label of the nearest training cycle, without using any feature, gives 0.01--0.15 times the error of the model evaluated correctly, as the bound of Section~\ref{sec:data} predicts. Per-cell min--max normalisation over the whole life, as in the PINN4SOH code, lowers the error on XJTU (0.80) and HUST (0.46), where all cells end at almost the same SOH and the normalised cycle index, the fraction of life already elapsed, nearly determines SOH; it raises the error on TJU (1.51) and MIT (2.76), where end-of-life SOH varies and the normalisation erases the absolute level of the features. Either way the number does not describe deployment, where the life of a cell is unknown. A whole-life outlier filter and a normaliser fitted with test cells have smaller effects, but are leaks as well. All further results use the leakage-free protocol.}
\newcommand{\LeakRowAbstract}{1.2--1.8}
\newcommand{\LeakMinmaxAbstract}{0.46--2.76}
\newcommand{\LeakMaxFactor}{2.8}
\newcommand{\BenchText}{PINN4SOH~\citep{wang2024} is the reference result on these four datasets. We reconstructed its 39 evaluation groups (one model per batch for XJTU and TJU, one per dataset for MIT and HUST with 23 and 20 fixed test cells, and few-shot runs with $K = 1, \dots, 4$ labelled cells for XJTU and HUST) from the \texttt{data path} lines of the training logs released with the code, and evaluate everything as the paper does: MAE per test cell, averaged over cells, runs and groups.

\paragraph{The released configuration.} The training logs record, for every epoch, the mean ``physics loss'' of the monotonicity term. In the first released version of the model this term is \path{relu(y2-y1).sum()}, computed from the labels of consecutive cycles only. In all 39 groups the logged value is one constant across all epochs of all ten runs, and in all of them it equals this expression evaluated on the training rows that the original data loader builds from our copy of the data (e.g.\ 0.043214 for the XJTU 2C batch). The term therefore adds a constant and no gradient: the published numbers come from a model trained with the data and PDE terms only. The match also shows that our reconstructed splits and data reproduce the training rows of every published group. Two later revisions replaced the term, first by \path{relu(u2-u1).sum()} and then by \path{relu((u2-u1)*(y1-y2)).sum()}, and the latest commit notes that the loss weights $\alpha$ and $\beta$ changed after a bug fix; the default weights in the released scripts no longer match those in the logs.

\paragraph{Re-running the original code.} We ran the released model and data-loading classes unchanged on the reconstructed splits, with ten seeded runs per group as in the paper and two versions of the model: the version that produced the published numbers, with the logged hyperparameters, and the current version with its default weights (Table~\ref{tab:bench}, Figure~\ref{fig:bench}). The original code needs pandas~1.x semantics; we used pandas~2.3, PyTorch~2 and CPUs, whereas the original runs used GPUs and older libraries. The re-run of the published version gives 1.04--1.19 times the published error on the four full datasets and 0.97--1.20 times in the few-shot settings. Per test cell, the re-run is significantly worse on 2 of the 4 full datasets. The standard deviation over runs is 4--34\,\% of the mean, so single runs differ by as much as these gaps. Since the training rows are identical, the remaining differences are hardware, library versions and random seeds; we could not reproduce the original environment and do not know which of them causes the systematic part. We therefore report both the published numbers and the re-run, and use the re-run for paired tests because it gives per-cell predictions under known conditions. 

\paragraph{Same splits, same preprocessing.} With the original preprocessing (whole-life $3\sigma$ filter on all columns including capacity, per-cell min--max normalisation of features and cycle index), PINN-hybrid in its fixed configuration ($\omega = \tfrac12$, $w = 10$, three seeds, no tuning on these splits) has 0.78--0.92 times the error of the re-run original code on the full datasets (significantly lower in 2 and higher in 0 of 4) and 0.31--0.67 times in the few-shot settings (lower in 4, higher in 0 of 8). Against the published numbers the ratios are 0.85--1.10 and 0.37--0.73. In the few-shot settings the unlabelled training cells enter the physics losses of our model, which is where the semi-supervised design pays off. The strongest model under these conditions, however, is not a PINN: a GBDT with the same preprocessing has 0.77--0.84 times the re-run error on the full datasets and 0.21--0.58 times in the few-shot settings, and is significantly lower than PINN-hybrid in 3 of the 12 settings (higher in 0; ratio 0.69--0.98). With per-cell min--max normalisation, the normalised cycle index is the fraction of life already elapsed, a feature that a tree exploits directly.

\paragraph{Without the leaks.} With our leakage-free preprocessing the same model has 0.63 (TJU), 0.23 (MIT), 1.65 (HUST) times the published error on the full datasets (per test cell significantly lower on 2 and higher on 1 of 4). The direction follows the leakage audit (Section~\ref{sec:leak}): the original preprocessing helps on datasets where the elapsed fraction of life determines SOH (HUST) and hurts where the absolute feature level matters (TJU, MIT). A comparison that changes the preprocessing therefore cannot separate the model from the protocol, and the published numbers on the few-shot settings, where knowing the elapsed fraction of life is almost knowing SOH, mostly reflect the normalisation.}
\newcommand{\BenchAbstract}{has 0.78--0.92 times the error of the re-run original code on the full datasets and 0.31--0.67 times in its few-shot settings}
\newcommand{\BenchDiscussion}{Three practical points follow. Reported numbers should come with per-cell errors and the exact splits, which made the comparison possible here; the logs released with PINN4SOH allowed us to reconstruct both. Released code should be the code that produced the numbers, or the differences should be listed: a constant loss term and changed weights alter the method that is being compared. And preprocessing that uses the whole life of a cell should be avoided, because its effect on the error depends on the dataset and can be large in either direction.}
